\begin{proof}
Strengthen the claim: for every derivation \(D\) with open assumptions
\(\Gamma\) and conclusion \(\chi\), and for every finite list
\(u_1,\ldots,u_r\) of pairwise distinct variables containing every variable
free in \(\Gamma\cup\set{\chi}\), the theory \(\ZF\) derives the displayed
guarded implication formed from \(\Gamma\), \(\chi\), and
\(u_1,\ldots,u_r\).  The lemma is the case in which \(D\) derives \(\psi\)
from \(\varphi_1,\ldots,\varphi_k\) and the list is \(v_1,\ldots,v_m\).
Induct on \(D\); throughout, argue inside \(\ZF\) under the quantifiers
\(\forall u_1\cdots\forall u_r\) and the guard hypotheses \(u_j\in L\).

Two bookkeeping devices recur.  First, \(\ZF\derives\exists u\,(u\in L)\):
Definition~\ref{def:p06-def-empty-base} gives
\[
  L_1=\operatorname{Def}(L_0)
     =\operatorname{Def}(\varnothing)=\{\varnothing\},
  \qquad \varnothing\in L_1\subseteq L.
\]
Thus \(\varnothing\) supplies the required witness directly.
Whenever the induction hypothesis supplies a claim
whose variable list contains a variable \(u\) absent from the target node,
instantiate the universally quantified guarded implication at this witness;
the vanished variable and its guard disappear.  Second, where the
capture-avoiding substitution of the calculus renames bound variables, the
relativizations of the formulas concerned are alpha-variants, because
\((\cdot)^{L}\) inserts the same parameter-free guard under alpha-variant
binders; the explicit \(\alpha\) rule of \(\ND_{\FOL}\) bridges them.

\emph{Assumptions, propositional rules, and classical rules.}  An
assumption leaf \(\chi\) maps to the relativized hypothesis \(\chi^{L}\) of
the guarded implication.  The operation \((\cdot)^{L}\) commutes with the
propositional connectives and \(\bottom^{L}=\bottom\), so every
propositional and classical rule instance maps to the same rule instance on
the relativized premises, with discharged assumptions mapped to discharged
relativized assumptions.

\emph{Quantifier rules.}  Every \(\Lang_{\in}\)-term is a variable, since
the signature has no function or constant symbols.  A \(\forall E\) step
from \(\forall x\,\chi\) to \(\chi[u/x]\) maps to a \(\forall E\) step on
\((\forall x\,\chi)^{L}=\forall x\,(x\in L\to\chi^{L})\), followed by
implication elimination against the guard \(u\in L\) supplied by the
antecedent.  An \(\exists I\) step from \(\chi[u/x]\) to
\(\exists x\,\chi\) maps to conjunction introduction with the guard
\(u\in L\) followed by \(\exists I\), concluding
\(\exists x\,(x\in L\land\chi^{L})\).  In both steps, if \(x\) is not free
in \(\chi\), the instantiating variable is absent from the formulas; use
the witness supplied by \(\exists u\,(u\in L)\) in place of a guard from
the antecedent.  A \(\forall I\) step with generalized variable \(x\) maps
to implication introduction discharging the hypothesis \(x\in L\),
followed by \(\forall I\); if \(x\) is not free in the premise, first
conjoin the hypothesis \(x\in L\) to the relativized premise and project
the premise back out, so that the discharged assumption occurs in its
derivation.  An \(\exists E\) step with eigenvariable \(y\) maps to
\(\exists E\) on the relativized existential premise; the minor derivation
assumes \(y\in L\land\chi^{L}[y/x]\), and conjunction eliminations recover
the guard \(y\in L\) and the relativized assumption.  The eigenvariable
conditions are unchanged, because the inserted guard \(x\in L\) is
parameter-free and adds no free variables.

\emph{Equality rules.}  Atomic formulas are unchanged by \((\cdot)^{L}\),
so an \(=I\) step concludes the unchanged atom.  Substitution of variables
for variables commutes with \((\cdot)^{L}\) up to the alpha-variant
bookkeeping above, so an \(=E\) step on \(\varphi\) maps to the \(=E\) step
on \(\varphi^{L}\).

\emph{Alpha rule.}  An explicit \(\alpha\) step maps to an \(\alpha\) step
by the second bookkeeping device.

At the root, the open relativized assumptions are
\(\varphi_1^{L},\ldots,\varphi_k^{L}\) and the guarded variables are
\(v_1,\ldots,v_m\); implication introduction and universal introduction in
the ambient \(\ZF\) derivation yield the displayed conclusion.
\end{proof}
